% =====================================================================
% DBA - Proyecto final - Informe técnico final
% Potencial solar en techos de los municipios PDET
% Integra las entregas 1 a 4 en un documento autocontenido para la UPME.
% =====================================================================
\documentclass[11pt,a4paper]{article}
\usepackage[utf8]{inputenc}
\usepackage[T1]{fontenc}
\IfFileExists{lmodern.sty}{\usepackage{lmodern}}{}
\usepackage[spanish,es-noshorthands,es-tabla]{babel}
\usepackage[margin=2.5cm]{geometry}
\usepackage{amsmath,amssymb}
\usepackage{booktabs,tabularx,array}
\usepackage{enumitem}
\usepackage[dvipsnames,table]{xcolor}
\usepackage{graphicx}
\usepackage{tikz}
\usetikzlibrary{calc,positioning,arrows.meta}
\usepackage{listings}
\usepackage{caption}
\usepackage{float}
\usepackage[hidelinks]{hyperref}
\usepackage{xurl}

% --- Colores --------------------------------------------------------
\definecolor{acento}{HTML}{2F6FDB}
\definecolor{gris}{HTML}{6B6B6B}
\hypersetup{colorlinks=true, linkcolor=acento, urlcolor=acento, citecolor=acento}

% --- Tablas ---------------------------------------------------------
\newcolumntype{L}{>{\raggedright\arraybackslash}X}
\renewcommand{\arraystretch}{1.25}

% --- Listas ---------------------------------------------------------
\setlist{itemsep=2pt, topsep=4pt}

% Más holgura entre palabras antes que salirse del margen con nombres en \texttt
\setlength{\emergencystretch}{3em}

% --- Código ---------------------------------------------------------
\lstdefinestyle{codigo}{
  basicstyle=\ttfamily\footnotesize,
  breaklines=true,
  columns=fullflexible,
  keepspaces=true,
  frame=single,
  framerule=0.3pt,
  rulecolor=\color{gray!50},
  backgroundcolor=\color{gray!6},
  xleftmargin=4pt, xrightmargin=4pt,
  showstringspaces=false,
  literate={á}{{\'a}}1 {é}{{\'e}}1 {í}{{\'i}}1 {ó}{{\'o}}1 {ú}{{\'u}}1
           {Á}{{\'A}}1 {É}{{\'E}}1 {Í}{{\'I}}1 {Ó}{{\'O}}1 {Ú}{{\'U}}1
           {ñ}{{\~n}}1 {Ñ}{{\~N}}1
}
\lstset{style=codigo}

% --- Atajos ---------------------------------------------------------
\newcommand{\campo}[1]{\texttt{#1}}
\newcommand{\ruta}[1]{\begingroup\urlstyle{tt}\url{#1}\endgroup}

% Figura de una captura o gráfica de img/.
%   \fig[ancho]{archivo}{alto relativo}{pie}{etiqueta}
\newcommand{\fig}[5][\textwidth]{%
  \begin{figure}[H]
    \centering
    \includegraphics[width=#1,height=#3\textheight,keepaspectratio]{img/#2}
    \caption{#4}
    \label{#5}
  \end{figure}}

% Las tablas generadas por el código (tablas/*.tex) vienen de la carpeta de la
% Entrega 4 y de src/analisis/potencial.py.
\newcommand{\tabla}[1]{\input{tablas/#1.tex}}

\begin{document}

% ------------------------------- Portada -----------------------------
\begin{titlepage}
  \centering
  {\large Pontificia Universidad Javeriana\par}
  {\large Departamento de Ingeniería de Sistemas\par}
  \vspace{0.8cm}
  {\large Administración de Bases de Datos (DBA) -- 2026-30\par}
  \vspace{2.2cm}
  {\Huge\bfseries Informe técnico final\par}
  \vspace{1cm}
  {\LARGE Potencial solar en techos\\ de los municipios PDET\par}
  \vspace{0.8cm}
  {\large Conteo de edificaciones y área de techos por municipio\\
  con MongoDB y dos fuentes abiertas de huellas de edificios\par}
  \vspace{0.6cm}
  {\large Preparado para la Oficina de Investigación e Información de la UPME\par}
  \vspace{2.2cm}
  {\large
    Maria Camila Beltran \\
    \vspace{0.5cm}
    Jorge Fortich \\
    \vspace{0.5cm}
    Jose Eduardo Valeriano \\
    \vspace{0.5cm}
    Samuel Osorio Rodríguez\par}
  \vspace{2.2cm}
  {\large Profesor: Andrés Oswaldo Calderón Romero, Ph.D.\par}
  \vspace{1.5cm}
  {\large Noviembre de 2026\par}
\end{titlepage}

% =====================================================================
% Resumen ejecutivo
% =====================================================================
\section*{Resumen ejecutivo}
\addcontentsline{toc}{section}{Resumen ejecutivo}

La Oficina de Investigación e Información de la UPME busca municipios PDET
candidatos para pilotos de energía solar. Este informe responde la pregunta
central del proyecto: \textbf{cuántas edificaciones hay y cuánta área de
techos tiene cada uno de los 170 municipios PDET}, según dos fuentes abiertas
de huellas de edificios, con una metodología reproducible construida sobre
MongoDB.

\textbf{Qué se hizo.} Se cargaron en MongoDB los límites oficiales de los 170
municipios (MGN 2025 del DANE) y las huellas de edificios de Google Open
Buildings v3 (2\,694\,786 edificios en PDET) y Microsoft Building Footprints
(1\,761\,048). Cada edificio se asignó a un único municipio por su centroide
con una consulta \campo{\$geoWithin} apoyada en un índice \campo{2dsphere}, y
las métricas por municipio se calcularon con un pipeline de agregación dentro
de la base. Todo el flujo, de la descarga a los mapas, corre con un solo
comando (\campo{make all}) y se reprodujo desde una base vacía.

\begin{table}[H]
  \centering\small
  \caption{Cifras principales (Google con confianza $\geq 0{,}65$).}
  \label{tab:resumen-ejecutivo}
  \begin{tabularx}{\textwidth}{@{}L r r@{}}
    \toprule
    & \textbf{Google} & \textbf{Microsoft} \\
    \midrule
    Municipios PDET (16 subregiones, 19 departamentos) & \multicolumn{2}{c}{170} \\
    Edificios asignados a un municipio & 2\,694\,785 & 1\,761\,048 \\
    Área total de techos & 222,2\,km\textsuperscript{2} & 221,1\,km\textsuperscript{2} \\
    Área media por edificio & 82\,m\textsuperscript{2} & 126\,m\textsuperscript{2} \\
    Municipios sin ningún edificio & 5 & 0 \\
    \midrule
    Área respaldada por las dos fuentes, $\sum \min(A_G, A_M)$ & \multicolumn{2}{c}{201,6\,km\textsuperscript{2} (20\,156\,ha)} \\
    Potencia instalable estimada (orden de magnitud) & \multicolumn{2}{c}{16\,GWp (rango 12 a 20\,GWp)} \\
    \bottomrule
  \end{tabularx}
\end{table}

\textbf{Hallazgos.}
\begin{itemize}
  \item Las dos fuentes coinciden en el total (222 y 221\,km\textsuperscript{2})
        y en el orden de los municipios (correlación de Spearman 0,94), aunque
        Google divide la misma superficie en 1,5 veces más huellas, más
        pequeñas. Por eso se compara y se ordena por \emph{área}, no por
        número de edificios.
  \item No coinciden en el Pacífico: en el Chocó y la costa nariñense Google
        detecta muy poco o nada (cinco municipios sin edificios). 18 de los
        23 municipios donde las fuentes están en fuerte desacuerdo son de esa
        región.
  \item El área de techos se concentra: los 15 primeros municipios suman el
        39\,\% del área conservadora de los 170.
  \item El resultado es estable: subir el umbral de confianza de Google a
        0,70 o 0,75 reduce el total, pero no cambia cuáles municipios
        encabezan la lista.
\end{itemize}

\textbf{Recomendaciones para la UPME.}
\begin{enumerate}
  \item Priorizar los pilotos en los municipios del top 15 por área
        conservadora (tabla~\ref{tab:ranking}), empezando por los que además
        tienen el área concentrada en un casco urbano: \textbf{Santa Marta},
        \textbf{Santander de Quilichao} y \textbf{Apartadó}.
  \item No descartar los municipios del Chocó y del Pacífico por su área
        baja: ahí la cifra depende de la fuente, y deben verificarse con
        imágenes recientes o en campo antes de decidir.
  \item Antes de dimensionar un piloto, refinar el área útil y cruzarla con
        la irradiación y la demanda local (sección~\ref{sec:recomendaciones}).
\end{enumerate}

\newpage


\tableofcontents
\newpage

% =====================================================================
% 1. Introducción
% =====================================================================
\section{Introducción}
\label{sec:intro}

\subsection{Contexto y objetivo}

La UPME evalúa la factibilidad de granjas solares en distintas regiones del
país y quiere priorizar los municipios designados como PDET (Programas de
Desarrollo con Enfoque Territorial), un foco de la política de desarrollo
posconflicto~\cite{pdet}. A mayor área de techos en un municipio, mayor
potencial para captar energía con paneles instalados sobre ellos. El
enunciado del proyecto~\cite{enunciado} pide una metodología reproducible
para \textbf{contar las edificaciones y sumar el área de sus techos en cada
municipio PDET}, comparando al menos dos fuentes abiertas de huellas de
edificios y usando una solución NoSQL.

El listado oficial (\campo{MunicipiosPDET.xlsx}) tiene 170 municipios en 16
subregiones PDET y 19 departamentos. Los límites municipales salen del Marco
Geoestadístico Nacional (MGN) del DANE~\cite{mgn}.

\subsection{Alcance}

\begin{itemize}
  \item \textbf{Territorio:} los 170 municipios PDET, con sus límites del MGN
        2025 (389\,182\,km\textsuperscript{2}).
  \item \textbf{Fuentes de edificios:} Google Open Buildings v3~\cite{google}
        y Microsoft Building Footprints~\cite{microsoft}, dos de las tres que
        propone el enunciado (la tercera es GlobalBuildingAtlas). La
        justificación está en la tabla~\ref{tab:decisiones}.
  \item \textbf{Métricas por municipio y fuente:} número de edificios, área
        total de techos, área media, densidad (edificios por
        km\textsuperscript{2}) y porcentaje del municipio cubierto por techos,
        con un análisis de sensibilidad al umbral de confianza de Google.
  \item \textbf{Fuera de alcance:} el diseño de los pilotos, la irradiación
        por municipio y la demanda eléctrica. La sección~\ref{sec:potencial}
        da solo un orden de magnitud de la potencia instalable.
\end{itemize}

El área de techo de cada edificio se aproxima con el área de su huella en
planta, que es lo que entregan las dos fuentes. Para cada municipio $m$ y
fuente $f$:

\begin{equation}
  A_{m,f} = \sum_{e \in E_f,\; \mathrm{centroide}(e) \in m} \mathrm{area}_{9377}(e)
  \label{eq:area}
\end{equation}

donde $E_f$ es el conjunto de edificios de la fuente $f$ y
$\mathrm{area}_{9377}(e)$ es el área de la huella calculada en la proyección
plana oficial de Colombia (EPSG:9377, en metros).

\subsection{Cómo se construyó y cómo leer este informe}

El proyecto se desarrolló en cuatro entregas semanales, cada una con su
código en el repositorio y su informe: diseño del esquema (semana 1), límites
PDET (semana 2), huellas de edificios (semana 3) y análisis (semana 4). Este
informe las integra en un solo documento:

\begin{itemize}
  \item Secciones~\ref{sec:arquitectura} y~\ref{sec:modelo}: arquitectura,
        decisiones tecnológicas y modelo de datos.
  \item Secciones~\ref{sec:municipios} a~\ref{sec:eda}: adquisición,
        integridad y carga de los datos, y su análisis exploratorio.
  \item Sección~\ref{sec:metodologia}: cómo se calcularon las métricas.
  \item Secciones~\ref{sec:resultados} y~\ref{sec:discusion}: resultados,
        discusión y recomendaciones.
  \item Sección~\ref{sec:reproduccion}: cómo reproducir todo.
  \item Anexo~\ref{sec:anexo}: capturas de evidencia de cada etapa.
\end{itemize}

% =====================================================================
% 2. Arquitectura y decisiones tecnológicas
% =====================================================================
\section{Arquitectura y decisiones tecnológicas}
\label{sec:arquitectura}

\subsection{Decisiones}

\begin{table}[H]
\centering\small
\caption{Decisiones tecnológicas del proyecto.}
\label{tab:decisiones}
\begin{tabularx}{\textwidth}{@{}p{2.6cm}p{3.6cm}L@{}}
\toprule
\textbf{Decisión} & \textbf{Elección} & \textbf{Justificación} \\
\midrule
Motor NoSQL & MongoDB 8 (modelo documental) &
  Guarda geometrías en GeoJSON de forma nativa, tiene índices \campo{2dsphere}
  y operadores \campo{\$geoWithin} y \campo{\$geoIntersects}~\cite{mongodb}.
  Su pipeline de agregación (\campo{\$group}, \campo{\$lookup},
  \campo{\$merge}) calcula los totales por municipio dentro de la base, y
  puede escalar horizontalmente con \emph{sharding}. \\
Fuente de edificios 1 & Google Open Buildings v3 &
  Cubre América Latina con imágenes de alta resolución (inferencia de mayo
  de 2023). Trae \campo{area\_in\_meters} y \campo{confidence} (0,65 a 1,0),
  lo que permite validar las áreas y medir la sensibilidad. Licencia
  CC BY 4.0 / ODbL. \\
Fuente de edificios 2 & Microsoft Building Footprints &
  Imágenes de 2014 a 2021, licencia ODbL. Particionado por país y
  \emph{quadkey}, así que se descarga solo lo necesario. Solo trae la
  geometría: el área se calcula. \\
Límites municipales & MGN 2025 del DANE, nivel municipio &
  Fuente oficial que exige el enunciado; se filtra con los 170 códigos del
  listado PDET. \\
Sistemas de referencia & WGS84 (EPSG:4326) para almacenar; MAGNA-SIRGAS
  Origen Nacional (EPSG:9377) para áreas &
  GeoJSON y \campo{2dsphere} exigen \campo{[longitud, latitud]} en WGS84. Las
  áreas se calculan en la proyección plana oficial de Colombia, en metros. \\
Herramientas & Python, GeoPandas, Shapely, PyMongo, DuckDB, Docker Compose,
  GitHub &
  GeoPandas valida y reproyecta; DuckDB lee y filtra los archivos grandes sin
  cargarlos completos; PyMongo carga por lotes; Docker Compose levanta la
  misma base en cualquier máquina. \\
\bottomrule
\end{tabularx}
\end{table}

\subsection{Flujo de datos}

El flujo tiene cuatro etapas encadenadas por un \campo{Makefile}
(figura~\ref{fig:flujo}). Cada etapa es una carpeta de \campo{src/} con un
archivo por tarea, lee sus parámetros de \campo{config.yaml} y deja un
registro en la colección \campo{bitacora\_cargas}. Las etapas se comunican por
archivos intermedios en \campo{data/processed/} y por las colecciones de
MongoDB, así que cada una se puede volver a correr por separado.

\begin{figure}[H]
\centering
\begin{tikzpicture}[
  font=\small,
  caja/.style={draw=acento, rounded corners=3pt, fill=acento!6, align=center,
               minimum height=1.25cm, text width=3.0cm},
  db/.style={draw=gris, rounded corners=3pt, fill=gray!8, align=center, font=\footnotesize,
             minimum height=0.9cm, text width=3.2cm},
  flecha/.style={-{Stealth[length=2.2mm]}, thick, gris},
  node distance=0.55cm and 0.55cm]
  \node[caja] (m) {\textbf{make municipios}\\ MGN 2025 + PDET\\ validación, carga, índice};
  \node[caja, right=of m] (e) {\textbf{make edificios}\\ Google y Microsoft\\ prefiltro, limpieza, carga};
  \node[caja, right=of e] (a) {\textbf{make analisis}\\ asignación, agregación,\\ sensibilidad};
  \node[caja, right=of a] (s) {\textbf{Salidas}\\ CSV, tablas, mapas,\\ ranking, potencial};
  \node[db, below=of m] (dm) {\campo{municipios}\\ \campo{fuentes}};
  \node[db, below=of e] (de) {\campo{edificios\_google}\\ \campo{edificios\_microsoft}};
  \node[db, below=of a] (da) {\campo{resultados}};
  \node[db, below=of s] (ds) {\ruta{data/processed/}\\ informe};
  \draw[flecha] (m) -- (e); \draw[flecha] (e) -- (a); \draw[flecha] (a) -- (s);
  \draw[flecha] (m) -- (dm); \draw[flecha] (e) -- (de); \draw[flecha] (a) -- (da);
  \draw[flecha] (s) -- (ds);
  \draw[flecha, dashed] (dm.east) -- (de.west);
  \draw[flecha, dashed] (de.east) -- (da.west);
  \draw[flecha, dashed] (da.east) -- (ds.west);
\end{tikzpicture}
\caption{Etapas del flujo (\campo{make all = make up municipios edificios
analisis}) y colecciones que produce cada una. Todas registran su ejecución en
\campo{bitacora\_cargas}.}
\label{fig:flujo}
\end{figure}

\subsection{Organización del repositorio y del trabajo}

\begin{lstlisting}
src/common/      conexión a MongoDB y lectura de config.yaml
src/municipios/  límites PDET: descarga, cruce, geometría, esquema, carga, índices
src/edificios/   huellas: descarga y limpieza por fuente, esquema, carga, índices
src/analisis/    asignación, agregación, sensibilidad, figuras, mapas, ranking
config.yaml      parámetros (CRS, umbrales, rutas, colecciones)
data/            datos crudos y procesados (solo se versionan las salidas pequeñas)
Entrega N/       informe de cada entrega; Informe Final/ este documento
\end{lstlisting}

El equipo trabajó con una rama y un \emph{pull request} por persona y por
tarea, repartidas por archivo para que las cuatro personas avanzaran en
paralelo. La tabla~\ref{tab:reparto} resume quién hizo qué.

\begin{table}[H]
\centering\small
\caption{Reparto del trabajo por semana.}
\label{tab:reparto}
\begin{tabularx}{\textwidth}{@{}p{2.4cm}LLLL@{}}
\toprule
\textbf{Semana} & \textbf{Jorge} & \textbf{Maria Camila} & \textbf{Samuel} & \textbf{Jose Eduardo} \\
\midrule
2. Límites PDET & Descarga del MGN, cruce y verificación & Geometría y exportación & Esquema, carga, índices y evidencia & Mapa de control y tamaño BSON \\
3. Edificios & Descarga y limpieza de Google & Descarga y limpieza de Microsoft & Esquema, carga, índices y \campo{make edificios} & EDA \\
4. Análisis & Asignación, \campo{explain} y validación & Mapas, ranking, \campo{run.py} e integración & Agregación, sensibilidad y exportación & Estilo de figuras y lectura de resultados \\
\bottomrule
\end{tabularx}
\end{table}

% =====================================================================
% 3. Modelo de datos
% =====================================================================
\section{Modelo de datos}
\label{sec:modelo}

\subsection{Diagrama E-R y decisiones de modelado}

El esquema vive en la base \campo{upme\_solar} y tiene seis colecciones. Las
centrales son los municipios y los edificios, que se relacionan por una
consulta espacial: cada edificio pertenece al municipio que contiene su
centroide (figura~\ref{fig:er}).

\fig{diagrama_er.png}{0.45}{Diagrama E-R del esquema en MongoDB (notación pata
de gallo). La línea azul es la relación espacial: \campo{cod\_mpio} se llena con
\campo{\$geoWithin}. $\blacklozenge$ marca un índice \campo{2dsphere}; los
campos indentados son subdocumentos embebidos.}{fig:er}

Cada relación es una referencia por \campo{\_id}: dos barras indican
exactamente uno, círculo y barra cero o uno, y círculo con pata de gallo cero
o muchos. Las decisiones principales fueron:

\begin{itemize}
  \item \textbf{Subregión PDET embebida en el municipio.} Es un dato pequeño,
        estable y que siempre se lee junto con el municipio.
  \item \textbf{Edificios referenciados, no embebidos.} Un municipio puede
        tener cientos de miles de edificios; embebidos superarían el límite de
        16\,MB por documento. Cada edificio guarda \campo{cod\_mpio}.
  \item \textbf{Una colección por fuente de edificios}, con el mismo esquema,
        para cargarlas, reindexarlas y compararlas por separado.
  \item \textbf{Centroide guardado aparte.} La asignación solo necesita un
        punto, e indexar un punto es más barato que indexar un polígono.
  \item \textbf{Resultados precalculados} en \campo{resultados}, para que los
        mapas y el informe no repitan la consulta espacial.
  \item \textbf{Trazabilidad:} la colección \campo{fuentes} guarda el origen y el
        SHA-256 de cada descarga, y \campo{bitacora\_cargas} cada ejecución
        con sus parámetros.
\end{itemize}

\campo{municipios}, \campo{fuentes}, las dos de edificios y \campo{resultados}
tienen un validador \campo{\$jsonSchema} derivado de este diccionario, así que MongoDB rechaza al insertar cualquier
documento que no lo cumpla (error 121; anexo~\ref{sec:anexo}).

\subsection{Diccionario de datos}

Todas las geometrías están en GeoJSON con \campo{[longitud, latitud]} en WGS84,
y todas las áreas en m\textsuperscript{2} calculados en EPSG:9377. Las tablas
reflejan el esquema final, con los ajustes hechos durante el proyecto.

\begin{table}[H]
\centering\small
\caption{Colección \campo{municipios} (170 documentos).}
\label{tab:dicc-mpio}
\begin{tabularx}{\textwidth}{@{}p{3cm}p{3.6cm}Lp{2.4cm}@{}}
\toprule
\textbf{Campo} & \textbf{Tipo} & \textbf{Descripción} & \textbf{Origen} \\
\midrule
\campo{\_id} & string & Código DANE a 5 dígitos (\campo{05045}) & MGN / Excel PDET \\
\campo{nombre} & string & Nombre del municipio & MGN \\
\campo{cod\_dpto}, \campo{departamento} & string & Código y nombre del departamento & MGN \\
\campo{subregion\_pdet} & objeto \{id: int, nombre\} & Subregión PDET (una de 16), embebida & Excel PDET \\
\campo{area\_km2} & double & Área del municipio & Calculado (EPSG:9377) \\
\campo{geometry} & GeoJSON Polygon o MultiPolygon & Límite municipal & MGN reproyectado \\
\campo{calidad} & objeto \{valida, reparada, n\_vertices\} & Resultado de la validación de la geometría & Calculado \\
\campo{fuente\_id} & string & Referencia a \campo{fuentes} (\campo{mgn}) & Calculado \\
\bottomrule
\end{tabularx}
\end{table}

\begin{table}[H]
\centering\small
\caption{Colecciones \campo{edificios\_google} y \campo{edificios\_microsoft}
(mismo esquema).}
\label{tab:dicc-edif}
\begin{tabularx}{\textwidth}{@{}p{2.8cm}p{3.1cm}Lp{2.4cm}@{}}
\toprule
\textbf{Campo} & \textbf{Tipo} & \textbf{Descripción} & \textbf{Origen} \\
\midrule
\campo{\_id} & ObjectId & Identificador interno & MongoDB \\
\campo{fuente\_id} & \campo{google} o \campo{microsoft} & Referencia a \campo{fuentes} & Calculado \\
\campo{geometry} & GeoJSON Polygon & Huella del edificio & Fuente \\
\campo{centroide} & GeoJSON Point & Centroide de la huella; base de la asignación & Calculado (EPSG:9377) \\
\campo{area\_m2} & double $\geq 0$ & Área de la huella & Calculado (EPSG:9377) \\
\campo{area\_m2\_fuente} & double o null & Área que reporta Google; null en Microsoft & Fuente \\
\campo{confianza} & double en $[0,1]$ o null & Confianza del modelo (0,65 a 1,0); null en Microsoft & Fuente \\
\campo{cod\_mpio} & string (5 dígitos) o null & Referencia a \campo{municipios}, llenada por la asignación & Calculado \\
\campo{geom\_hash} & string (índice único) & Hash de la geometría normalizada, contra duplicados & Calculado \\
\campo{lote\_id} & ObjectId & Carga que trajo el registro (\campo{bitacora\_cargas}) & Calculado \\
\bottomrule
\end{tabularx}
\end{table}

\begin{table}[H]
\centering\small
\caption{Colección \campo{resultados} (680 documentos: 170 municipios por
Google con tres umbrales y Microsoft).}
\label{tab:dicc-res}
\begin{tabularx}{\textwidth}{@{}p{3.2cm}p{3.6cm}L@{}}
\toprule
\textbf{Campo} & \textbf{Tipo} & \textbf{Descripción} \\
\midrule
\campo{\_id} & objeto \{cod\_mpio, fuente\_id, umbral\} & Un documento por municipio, fuente y umbral (\campo{umbral: null} en Microsoft) \\
\campo{n\_edificios} & int o long & Edificios cuyo centroide cae en el municipio \\
\campo{area\_total\_m2} & double & Suma de \campo{area\_m2} \\
\campo{area\_media\_m2} & double o null & Área promedio (null si no hay edificios) \\
\campo{densidad\_km2} & double & Edificios por km\textsuperscript{2} del municipio \\
\campo{pct\_techos} & double en $[0, 100]$ & Área de techos sobre el área del municipio, en \% \\
\campo{parametros} & objeto & Umbral, regla de asignación, CRS de las áreas y commit del código \\
\campo{fecha\_calculo} & date & Cuándo se calculó \\
\bottomrule
\end{tabularx}
\end{table}

\begin{table}[H]
\centering\small
\caption{Colecciones \campo{fuentes} y \campo{bitacora\_cargas}.}
\label{tab:dicc-fuentes}
\begin{tabularx}{\textwidth}{@{}p{3.8cm}p{2.4cm}L@{}}
\toprule
\textbf{Campo} & \textbf{Tipo} & \textbf{Descripción} \\
\midrule
\multicolumn{3}{@{}l}{\textbf{\campo{fuentes}} (una por dataset: \campo{mgn}, \campo{google}, \campo{microsoft})} \\
\campo{\_id}, \campo{nombre}, \campo{version}, \campo{licencia} & string & Identificación del dataset \\
\campo{url\_descarga}, \campo{crs\_origen} & string & Origen exacto y CRS en que llegan los datos \\
\campo{fecha\_descarga} & date & Cuándo se descargó \\
\campo{checksum\_sha256} & string & SHA-256 de la descarga \\
\campo{periodo\_imagenes} & string o null & Fechas de las imágenes (no aplica al MGN) \\
\campo{n\_registros} & long o null & Registros cargados; lo completa la carga \\
\midrule
\multicolumn{3}{@{}l}{\textbf{\campo{bitacora\_cargas}} (una por ejecución)} \\
\campo{fuente\_id}, \campo{etapa} & string & Fuente y etapa (\campo{descarga}, \campo{limpieza}, \campo{carga}, \campo{indices}, \campo{analisis}) \\
\campo{inicio}, \campo{fin} & date & Duración de la etapa \\
\campo{registros\_ok}, \campo{registros\_error} & long & Resultado de la etapa \\
\campo{parametros} & objeto & Copia de la configuración usada \\
\bottomrule
\end{tabularx}
\end{table}

\textbf{Cambios respecto al diseño inicial.} En \campo{fuentes},
\campo{periodo\_imagenes} y \campo{n\_registros} pasaron a aceptar \campo{null}
y se agregó \campo{checksum\_sha256}; exigirlos desde el primer
\emph{insert} bloqueaba el flujo, porque se completan en otra etapa. En
\campo{resultados}, la llave pasó de \campo{\{cod\_mpio, fuente\_id\}} a
\campo{\{cod\_mpio, fuente\_id, umbral\}} para guardar la sensibilidad en la
misma colección.

\subsection{Índices}

\begin{table}[H]
\centering\small
\caption{Índices del esquema.}
\label{tab:indices}
\begin{tabularx}{\textwidth}{@{}p{3.2cm}p{4.4cm}L@{}}
\toprule
\textbf{Colección} & \textbf{Índice} & \textbf{Para qué} \\
\midrule
\campo{municipios} & \campo{2dsphere} sobre \campo{geometry} & Consultas espaciales sobre los límites \\
\campo{municipios} & \campo{subregion\_pdet.id} & Resúmenes por subregión \\
\campo{edificios\_*} & \campo{2dsphere} sobre \campo{centroide} & Asignación con \campo{\$geoWithin} (consulta principal) \\
\campo{edificios\_*} & \campo{2dsphere} sobre \campo{geometry} & \campo{\$geoIntersects} sobre edificios en los límites \\
\campo{edificios\_*} & \campo{cod\_mpio} & Agregación por municipio \\
\campo{edificios\_*} & \campo{geom\_hash} (único) & Impedir duplicados \\
\campo{resultados} & \campo{\_id} compuesto & Búsqueda por municipio, fuente y umbral \\
\bottomrule
\end{tabularx}
\end{table}

% =====================================================================
% 4. Límites PDET (semana 2)
% =====================================================================
\section{Límites de los municipios PDET}
\label{sec:municipios}

\campo{make municipios} (\campo{src/municipios/run.py}) descarga el MGN 2025,
lo cruza con el listado PDET, valida y reproyecta las geometrías, y las carga
e indexa en MongoDB.

\subsection{Adquisición y verificación}
\label{sec:mgn-descarga}

\textbf{Descarga.} El geoportal del DANE genera la descarga desde una página
web, sin un enlace directo y estable, así que \campo{descarga.py} acepta el
ZIP puesto a mano en \ruta{data/raw/mgn_2025/} o una URL directa opcional
(\campo{MGN\_URL\_DIRECTA} en \campo{.env}). Se usó
\campo{MGN2025\_MPIO\_GRAFICO.zip} (nivel municipio, Shapefile). El script
calcula el SHA-256 del ZIP, lo descomprime y registra la fuente en
\campo{fuentes} (\campo{\_id="mgn"}) con URL, versión, CRS leído del
\campo{.prj}, fecha y checksum:

\begin{lstlisting}
SHA-256: a3c01393059be8f2947de5e0a828faa4c7a43119819a667a8124854340bc141f
\end{lstlisting}

La misma suma se obtuvo en la corrida desde base vacía de este informe
(sección~\ref{sec:reproduccion}), lo que confirma que el insumo no cambió.

\textbf{Cruce con el listado PDET.} \campo{cruce.py} lee
\campo{MunicipiosPDET.xlsx} y la capa municipal del MGN y cruza por código
DANE. Los códigos del Excel vienen como números y los de Antioquia pierden el
cero inicial (5045 en vez de 05045), así que se rellenan a cinco dígitos antes
del cruce. \campo{verificacion.py} es un control independiente: vuelve a
contar desde el Excel y la capa cruzada.

\begin{table}[H]
  \centering\small
  \caption{Control de cobertura del cruce (\campo{verificacion.py}).}
  \label{tab:verificacion}
  \begin{tabularx}{0.85\textwidth}{L r r c}
    \toprule
    Verificación & Esperado & Obtenido & OK \\
    \midrule
    Municipios en el Excel (códigos únicos) & 170 & 170 & \checkmark \\
    Códigos repetidos en el Excel           & 0   & 0   & \checkmark \\
    Municipios que cruzaron con el MGN      & 170 & 170 & \checkmark \\
    Subregiones PDET en el resultado        & 16  & 16  & \checkmark \\
    Departamentos en el resultado           & 19  & 19  & \checkmark \\
    \bottomrule
  \end{tabularx}
\end{table}

\subsection{Integridad y formato de las geometrías}
\label{sec:mgn-geometria}

MongoDB rechaza un polígono inválido al insertarlo en una colección con
índice \campo{2dsphere}, así que la validez se revisa antes de cargar.
\campo{geometria.py}:

\begin{itemize}
  \item Revisa \campo{is\_valid} (GEOS, vía Shapely) y guarda el motivo de
        las inválidas; las repara con \campo{make\_valid} conservando solo las
        partes poligonales, para que el resultado sea siempre
        \campo{Polygon} o \campo{MultiPolygon}.
  \item Calcula \campo{area\_km2} en EPSG:9377 y reproyecta a EPSG:4326 con
        orden \campo{[lon, lat]}.
  \item Ajusta las coordenadas a una rejilla de $10^{-6}$ grados (unos
        10\,cm), elimina vértices repetidos y orienta los anillos según
        RFC~7946.
  \item Controla que todas las coordenadas caigan en el rectángulo de
        Colombia (detecta lon/lat invertidas) y vuelve a revisar la validez
        después de reproyectar.
\end{itemize}

Las 170 geometrías llegaron válidas y ninguna necesitó reparación. Como
verificación independiente, el área calculada se comparó con el área oficial
del MGN (\campo{mpio\_narea}): la diferencia máxima fue de
$2{,}2\times10^{-5}$\,\% y el total coincide (389\,182\,km\textsuperscript{2}).

\begin{table}[H]
  \centering\small
  \caption{Resultados de la validación y exportación de los municipios.}
  \label{tab:geometria}
  \begin{tabularx}{0.85\textwidth}{L r}
    \toprule
    Indicador & Valor \\
    \midrule
    Municipios procesados / válidos de origen / reparados & 170 / 170 / 0 \\
    Área total de los municipios PDET  & 389\,182\,km\textsuperscript{2} \\
    Municipio más grande / más pequeño & Solano, 42\,234\,km\textsuperscript{2} / Chalán, 84\,km\textsuperscript{2} \\
    Vértices en total (máximo: Solano) & 1\,551\,994 (51\,568) \\
    Documento más pesado (Solano)      & 1,74\,MB (10,9\,\% del límite de 16\,MB) \\
    Tamaño de \campo{municipios\_pdet.geojson} & 37,1\,MB \\
    \bottomrule
  \end{tabularx}
\end{table}

\campo{exportar.py} escribe \ruta{data/processed/municipios_pdet.geojson} con
un \campo{Feature} por municipio y los campos del diccionario en
\campo{properties}; valida que no haya códigos repetidos y que todos tengan
cinco dígitos. Este archivo se versiona, para que nadie dependa de volver a
descargar el MGN. \campo{calidad.py} mide el tamaño BSON exacto de cada
documento con \campo{bson.encode()} y falla si alguno supera el límite; el
mayor ocupa el 10,9\,\% del límite, así que ninguno tuvo que simplificarse.

\subsection{Carga, índices y evidencia}
\label{sec:mgn-carga}

\campo{schema.py} crea \campo{municipios} y \campo{fuentes} con su validador.
\campo{carga.py} usa \campo{bulk\_write} con un \campo{UpdateOne(...,
upsert=True)} por municipio, así que volver a correrlo no duplica documentos;
al terminar completa \campo{fuentes.mgn.n\_registros} y registra la ejecución
en \campo{bitacora\_cargas}. \campo{indices.py} crea el \campo{2dsphere} sobre
\campo{geometry} y el índice de \campo{subregion\_pdet.id}.

\campo{evidencia.py} prueba que el índice se \emph{usa}: toma un punto interior
de un municipio (Amalfi), corre \campo{\$geoIntersects} y exige que el plan de
\campo{explain()} tenga \campo{IXSCAN} y no \campo{COLLSCAN}, y que la consulta
devuelva exactamente ese municipio.

\begin{lstlisting}
[evidencia] etapas del plan: ['FETCH', 'IXSCAN'] (se espera IXSCAN, no COLLSCAN)
[evidencia] punto de control: AMALFI (05031) en [-74.95147, 7.02048]
[evidencia] $geoIntersects devolvió: [{'_id': '05031', 'nombre': 'AMALFI'}]
\end{lstlisting}

\fig[0.8\textwidth]{cap_mapa.png}{0.4}{Mapa de control: los 170 municipios PDET
coloreados por sus 16 subregiones, sin huecos ni geometrías sueltas.}{fig:mapa-control}

% =====================================================================
% 5. Huellas de edificios (semana 3)
% =====================================================================
\section{Huellas de edificios}
\label{sec:edificios}

\campo{make edificios} (\campo{src/edificios/run.py}) procesa cada fuente
con el mismo esquema: descarga con SHA-256, prefiltro a PDET por la regla del
centroide, limpieza compartida (\campo{comun.py}) hacia un Parquet con el
mismo contrato de columnas, y carga e índices en MongoDB.

\begin{table}[H]
  \centering\small
  \caption{Adquisición y limpieza de las dos fuentes.}
  \label{tab:edif-resumen}
  \begin{tabularx}{\textwidth}{@{}L r r@{}}
    \toprule
    & \textbf{Google Open Buildings v3} & \textbf{Microsoft Building Footprints} \\
    \midrule
    Archivos descargados & 8 celdas S2 (3,11\,GB) & 125 \emph{quadkeys} (373,4\,MB) \\
    Edificios leídos & 16\,512\,141 (en Colombia) & 4\,595\,615 \\
    Edificios en PDET (en la colección) & 2\,694\,786 & 1\,761\,048 \\
    Polígonos inválidos reparados (GEOS) & 0 & 28 \\
    Reparados en la carga (validación esférica) & 0 & 29 \\
    Duplicados por \campo{geom\_hash} & 0 & 0 \\
    Licencia & CC BY 4.0 / ODbL & ODbL \\
    \bottomrule
  \end{tabularx}
\end{table}

\subsection{Google Open Buildings: descarga y prefiltro}
\label{sec:google}

Google publica sus edificios en CSV comprimidos, uno por celda S2 de nivel~4.
Cada edificio trae su polígono en WKT, el centroide, el área que calcula
Google (\campo{area\_in\_meters}) y la confianza del modelo
(\campo{confidence}, 0,65 a 1,0).

\textbf{Descarga.} \campo{descarga\_google.py} cruza la tabla de celdas que
publica Google con la unión de los 170 municipios: quedan 8 celdas (3,11\,GB)
de las 10 que tocan Colombia. Cada celda se baja en \emph{streaming}, se
reanuda si se corta, se verifica su tamaño y se calcula su SHA-256; el SHA-256
de la fuente (hash de la lista ordenada de pares \campo{token:sha256}) se
guarda en \campo{fuentes} y dio el mismo valor en dos máquinas distintas.

\textbf{Prefiltro con DuckDB.} Las 8 celdas tienen 31 millones de filas en el
rectángulo de Colombia, así que no se cargan completas: DuckDB lee cada CSV
comprimido en tandas de unas 100\,000 filas y descarta en SQL lo que cae fuera
de los rectángulos PDET. Luego se prueba con un índice espacial de Shapely si
el centroide cae dentro de algún municipio PDET. De 31\,370\,317 filas en el
rectángulo, 16\,512\,141 están dentro de Colombia y 2\,694\,786 en los
municipios PDET (16,3\,\%). Una corrida de control con las 10 celdas dio
exactamente los mismos 2\,694\,786 edificios.

\fig{cap_google_prefiltro.png}{0.28}{Prefiltro de Google: acumulados en
Colombia y tras el prefiltro por cada celda.}{fig:g-prefiltro}

\subsection{Microsoft Building Footprints: descarga y prefiltro}
\label{sec:microsoft}

Microsoft publica sus huellas globales (versión 2026-02-03) como GeoJSON por
líneas comprimido, en particiones por \emph{quadkey} de nivel 9; solo traen el
polígono. \campo{descarga\_microsoft.py} lee el índice público
\campo{dataset-links.csv}, toma las 232 particiones de Colombia, se queda con
las 125 que intersectan la unión de los municipios PDET y las baja en paralelo
(8 hilos), con un \campo{manifest.csv} que guarda el SHA-256 de cada archivo y
uno global. \campo{limpieza\_microsoft.py} procesa cada archivo en un proceso
aparte (8 en paralelo) y conserva solo los edificios cuyo centroide cae en
PDET: de 4\,595\,615 quedan 1\,761\,048.

\fig{cap_ms_limpieza.png}{0.25}{Tabla de control de la limpieza de
Microsoft.}{fig:ms-limpieza}

\subsection{Limpieza común y contrato del Parquet}
\label{sec:limpieza}

Las dos fuentes pasan por la misma limpieza, para que los dos Parquet
intermedios tengan las mismas columnas:

\begin{enumerate}[leftmargin=*]
  \item Se construyen las geometrías con Shapely y se descartan las ilegibles.
  \item \textbf{Validación y reparación:} los polígonos inválidos se reparan
        con \campo{make\_valid}; como el contrato pide \campo{Polygon}, se
        conserva la parte poligonal mayor.
  \item \textbf{Área y centroide:} se proyecta a EPSG:9377 para calcular
        \campo{area\_m2} y el centroide; la geometría se guarda en EPSG:4326.
  \item \textbf{\campo{geom\_hash}:} hash de la geometría normalizada y
        redondeada, para que dos copias del mismo edificio den el mismo valor;
        se descartan los repetidos (un edificio puede venir en dos archivos).
  \item DuckDB une las partes en un único Parquet comprimido (zstd) con
        metadatos GeoParquet.
\end{enumerate}

\begin{table}[H]
\centering\small
\caption{Contrato de columnas de los Parquet intermedios.}
\label{tab:contrato}
\begin{tabularx}{\textwidth}{@{}p{3.6cm}p{2.2cm}L@{}}
\toprule
\textbf{Columna} & \textbf{Tipo} & \textbf{Contenido} \\
\midrule
\campo{fuente\_id} & texto & \campo{google} o \campo{microsoft} \\
\campo{geometry} & WKB & \campo{Polygon} en EPSG:4326 \\
\campo{centroide\_lon}, \campo{centroide\_lat} & double & Centroide de la huella, en EPSG:4326 \\
\campo{area\_m2} & double & Área en EPSG:9377, en m\textsuperscript{2} \\
\campo{area\_m2\_fuente}, \campo{confianza} & double o null & Área y confianza de Google; null en Microsoft \\
\campo{geom\_hash} & texto & Hash de la geometría normalizada \\
\bottomrule
\end{tabularx}
\end{table}

\subsection{Carga eficiente en MongoDB}
\label{sec:carga}

Con 4,5 millones de documentos, la eficiencia de la carga depende de no
cargar todo a memoria y de no mantener índices mientras se inserta.
\campo{carga.py}:

\begin{itemize}
  \item Lee el Parquet con \campo{pyarrow} en \textbf{lotes de 5\,000 filas}
        (\campo{config.yaml}); la geometría pasa de WKB a GeoJSON con Shapely.
  \item Inserta con \campo{insert\_many(ordered=False)}: un documento
        rechazado no frena el resto del lote.
  \item \textbf{Crea los índices después de cargar}, porque mantenerlos
        durante la inserción es más lento. La excepción es el índice de
        \campo{geometry} (sección~\ref{sec:s2}).
  \item Es \textbf{idempotente}: borra los documentos previos de la fuente
        antes de cargar; correrla dos veces deja el mismo conteo.
  \item Marca cada edificio con un \campo{lote\_id}, que es el \campo{\_id} de
        su fila en \campo{bitacora\_cargas}, con inicio, fin, registros y
        velocidad.
\end{itemize}

\begin{table}[H]
\centering\small
\caption{Carga a MongoDB (MongoDB 8 en Docker local, un solo equipo; los
tiempos dependen de la máquina).}
\label{tab:carga}
\begin{tabularx}{\textwidth}{@{}lrrrrL@{}}
\toprule
\textbf{Fuente} & \textbf{Documentos} & \textbf{Rechazados} & \textbf{Reparados (S2)} & \textbf{Tiempo (s)} & \textbf{Velocidad} \\
\midrule
Google    & 2\,694\,786 & 0 & 0  & 56,6 & 47\,621 docs/s \\
Microsoft & 1\,761\,048 & 0 & 29 & 37,7 & 46\,753 docs/s \\
\bottomrule
\end{tabularx}
\end{table}

\begin{table}[H]
\centering\small
\caption{Tiempo de creación de los índices de edificios.}
\label{tab:indices-edif}
\begin{tabularx}{\textwidth}{@{}p{3.6cm}Lrr@{}}
\toprule
\textbf{Índice} & \textbf{Para qué} & \textbf{Google (s)} & \textbf{Microsoft (s)} \\
\midrule
\campo{2dsphere} en \campo{centroide} & \campo{\$geoWithin} de la asignación & 6,9 & 4,4 \\
\campo{2dsphere} en \campo{geometry} & \campo{\$geoIntersects} en límites & \multicolumn{2}{c}{creado antes de cargar} \\
\campo{cod\_mpio} & agregación por municipio & 2,2 & 1,4 \\
\campo{geom\_hash} (único) & impedir duplicados & 3,4 & 2,3 \\
\bottomrule
\end{tabularx}
\end{table}

En disco, \campo{edificios\_google} ocupa 1\,279\,MB de datos y 245\,MB de
índices, y \campo{edificios\_microsoft} 835\,MB y 175\,MB. El índice único
sobre \campo{geom\_hash} es una segunda barrera contra duplicados: si alguno
llegara a la colección, la creación del índice fallaría.

\subsection{Un hallazgo: polígonos válidos para GEOS que MongoDB rechaza}
\label{sec:s2}

Con todos los índices creados al final, el \campo{2dsphere} sobre
\campo{geometry} de Microsoft falló: MongoDB construye sus claves con geometría
esférica (biblioteca S2) y encontró un polígono con aristas que se cruzan,
aunque GEOS, que trabaja en el plano, lo había dado por válido. Con una sola
geometría así, la construcción de todo el índice aborta. En una muestra de
300\,000 edificios hubo 5 casos, unos 2 por millón.

La solución fue crear ese índice \emph{antes} de cargar: así MongoDB rechaza
solo los documentos afectados y la carga sigue. Cada rechazado se repara
redondeando las coordenadas a $10^{-6}$ grados (unos 11\,cm) con
\campo{shapely.set\_precision} y se reintenta; \campo{buffer(0)} y
\campo{simplify} no resolvieron ninguno de los casos de prueba. En la carga
completa, Microsoft tuvo 29 reparados y Google ninguno, y no se descartó ningún
edificio.

\subsection{Evidencia de que los índices se usan}
\label{sec:evidencia-edif}

\campo{evidencia.py} corre la consulta de asignación sobre el polígono de
Caldono (19137) con el índice y forzando el recorrido completo
(\campo{hint(\{\$natural: 1\})}).

\begin{table}[H]
\centering\small
\caption{Planes de ejecución en Caldono (una consulta por caso; tiempos
indicativos).}
\label{tab:evidencia-edif}
\begin{tabularx}{\textwidth}{@{}llLrr@{}}
\toprule
\textbf{Fuente} & \textbf{Plan} & \textbf{Consulta} & \textbf{Resultados} & \textbf{Tiempo (s)} \\
\midrule
Google    & IXSCAN   & \campo{\$geoWithin} sobre \campo{centroide}   & 15\,303 & 0,44 \\
Google    & IXSCAN   & \campo{\$geoIntersects} sobre \campo{geometry} & 15\,304 & 0,37 \\
Google    & COLLSCAN & \campo{\$geoWithin} sin índice                 & 15\,303 & 2,93 \\
Microsoft & IXSCAN   & \campo{\$geoWithin} sobre \campo{centroide}   & 12\,547 & 0,32 \\
Microsoft & IXSCAN   & \campo{\$geoIntersects} sobre \campo{geometry} & 12\,548 & 0,28 \\
Microsoft & COLLSCAN & \campo{\$geoWithin} sin índice                 & 12\,547 & 1,93 \\
\bottomrule
\end{tabularx}
\end{table}

La diferencia de un edificio entre \campo{\$geoWithin} (centroide dentro) y
\campo{\$geoIntersects} (huella que toca el municipio) es un edificio que cruza
el límite: es la razón de usar el centroide, para que cada edificio pertenezca
a un solo municipio.

% =====================================================================
% 6. Análisis exploratorio
% =====================================================================
\section{Análisis exploratorio de las huellas}
\label{sec:eda}

El EDA se hizo sobre los dos Parquet completos, sin muestreo
(\ruta{notebooks/eda_edificios.ipynb}); las figuras se regeneran con
\campo{src/analisis/figuras\_eda.py}. Sus cifras coinciden con los documentos
cargados en MongoDB.

\subsection{Totales y tamaño de las huellas}

\begin{table}[H]
\centering\small
\caption{Totales por fuente sobre los 170 municipios PDET (áreas en EPSG:9377).}
\label{tab:eda-totales}
\begin{tabularx}{\textwidth}{@{}Lrrrr@{}}
\toprule
\textbf{Fuente} & \textbf{Edificios} & \textbf{Área total (km\textsuperscript{2})} & \textbf{Área media (m\textsuperscript{2})} & \textbf{Mediana (m\textsuperscript{2})} \\
\midrule
Google Open Buildings & 2\,694\,786 & 222,17 & 82,45 & 57,87 \\
Microsoft Building Footprints & 1\,761\,048 & 221,13 & 125,57 & 76,99 \\
\bottomrule
\end{tabularx}
\end{table}

Google tiene 1,53 veces más edificios que Microsoft, pero más pequeños: su área
media es un tercio menor y su mediana un 25\,\% menor. Aun así, el área total
es casi igual (diferencia de 1,04\,km\textsuperscript{2}, menos del 0,5\,\%):
Google reparte la misma superficie en más huellas. En Microsoft la cola de
edificios grandes es más larga (p99 de 912\,m\textsuperscript{2} frente a
445\,m\textsuperscript{2}).

\fig{eda_confianza_area.png}{0.28}{Izquierda: distribución de áreas por fuente
(escala logarítmica, medianas marcadas). Derecha: distribución de la confianza
de Google, con los umbrales 0,70 y 0,75 señalados.}{fig:eda-confianza}

La confianza de Google va de 0,65 a 0,98 (media y mediana 0,78). Ningún
registro queda por debajo de 0,65, porque Google solo publica detecciones por
encima de ese valor; por eso 0,65 es el escenario base y 0,70 y 0,75 se usan
para medir la sensibilidad (sección~\ref{sec:sensibilidad}).

\subsection{Validación del área calculada}

Para Google se comparó \campo{area\_m2} (calculada en EPSG:9377) con
\campo{area\_m2\_fuente} (la de Google). La correlación es prácticamente
perfecta y la diferencia relativa tiene mediana 0 y percentiles 5 y 95 de
$\pm$1\,\%; las diferencias mayores son huellas muy pequeñas, donde un redondeo
pesa mucho. El cálculo de área no tiene sesgos sistemáticos.

\fig{eda_comparacion.png}{0.26}{Izquierda: área calculada frente a la de Google
(recta 1:1). Centro: diferencia relativa. Derecha: área total por fuente.}{fig:eda-comparacion}

\subsection{Distribución por subregión y cobertura}

Las dos fuentes concentran sus edificios en las mismas subregiones, pero en
distinta proporción: la Sierra Nevada - Perijá tiene 645\,009 edificios en
Google y 295\,383 en Microsoft. El Chocó es la subregión con menos edificios en
Google (20\,722) y no en Microsoft (37\,692): la primera señal de que las
fuentes no cubren el territorio igual.

\fig{eda_subregion.png}{0.3}{Edificios por subregión PDET y fuente.}{fig:eda-subregion}

El mapa de densidad de centroides (figura~\ref{fig:eda-mapa}) muestra dónde
están los huecos: cinco municipios no tienen ningún edificio en Google
(Murindó y Vigía del Fuerte en Antioquia; Medio San Juan, Nóvita y Bojayá en el
Chocó), mientras que Microsoft tiene entre 1\,331 y 2\,004 en cada uno. Otros
tienen muy pocos en Google, como Sipí (1), López de Micay (24) y Condoto (239).

\fig{eda_mapa_densidad.png}{0.4}{Densidad de centroides por fuente sobre los
municipios PDET (escala logarítmica común). En rojo, los cinco municipios sin
edificios de Google.}{fig:eda-mapa}

Las figuras de esta sección son la versión corregida de las del EDA original,
rehechas con un estilo común para todo el informe (textos en español, escalas
logarítmicas reales, miles separados y la misma escala de color en los dos
paneles), atendiendo la retroalimentación recibida sobre las gráficas.

% =====================================================================
% 7. Metodología del cálculo (semana 4)
% =====================================================================
\section{Metodología del cálculo}
\label{sec:metodologia}

Las métricas de la ecuación~(\ref{eq:area}) se calculan en dos pasos dentro
de MongoDB: primero se asigna cada edificio a un municipio
(\campo{asignacion.py}) y después se agrega por municipio, fuente y umbral
(\campo{agregacion.py}). Una validación independiente en GeoPandas comprueba
la asignación.

\subsection{Asignación por centroide}
\label{sec:asignacion}

\textbf{Regla.} Un edificio pertenece al municipio que contiene su centroide.
Así un techo que cruza el límite entre dos municipios se cuenta una sola vez,
y basta un índice sobre puntos, más barato que uno sobre polígonos.

\textbf{Procedimiento.} Para cada uno de los 170 municipios se ejecuta un
\campo{updateMany} con \campo{\$geoWithin} sobre \campo{centroide}, usando la
geometría del municipio como ventana:

\begin{lstlisting}[language=Python]
coleccion.update_many(
    {"cod_mpio": None,
     "centroide": {"$geoWithin": {"$geometry": municipio["geometry"]}}},
    {"$set": {"cod_mpio": municipio["_id"]}})
\end{lstlisting}

El filtro \campo{cod\_mpio: null} hace que el paso sea reanudable y que un
edificio no se reasigne. La ejecución queda en \campo{bitacora\_cargas}.

\tabla{asignacion_resumen}

En Microsoft todos los edificios quedaron asignados; en Google, todos menos
uno. Los tiempos de la tabla~\ref{tab:asignacion} son de una corrida en otra
máquina; en la corrida desde base vacía de la sección~\ref{sec:reproduccion}
la asignación tomó 101,5\,s en Google y 68,1\,s en Microsoft, con exactamente
las mismas cifras.

\textbf{El edificio sin municipio.} Su centroide (-72,97638; 9,15139) está a
0,011\,m del límite de Convención (54206), prácticamente sobre la frontera. En
un caso así, la diferencia entre las aristas geodésicas que usa
\campo{\$geoWithin} sobre la esfera y las aristas planas de la geometría puede
dejar el punto fuera. Es un edificio de 2,7 millones (0,00004\,\%), sin efecto
en los resultados.

\textbf{Plan de ejecución.} \campo{evidencia.py} corre
\campo{explain("executionStats")} sobre la consulta de asignación en Caldono
(19137), con el índice \campo{centroide\_2dsphere} y forzando
\campo{COLLSCAN} con \campo{hint(\{\$natural: 1\})}.

\tabla{evidencia_geowithin}

Con el índice, la consulta de Google examinó 25\,071 documentos en 730\,ms; sin
él recorrió los 2\,694\,786 en 7\,310\,ms: 107 veces más documentos y 10 veces
más tiempo. En Microsoft, 91 veces más documentos y 6 veces más tiempo. Las dos
variantes devuelven los mismos edificios. El índice examina más documentos de
los que devuelve porque primero filtra por celdas S2 y luego verifica la
geometría exacta.

\fig{cap_explain_geowithin.png}{0.3}{\campo{explain("executionStats")} con
\campo{IXSCAN} y con \campo{COLLSCAN}.}{fig:explain}

\subsection{Validación independiente con GeoPandas}
\label{sec:validacion}

\campo{validacion.py} repite la asignación fuera de MongoDB en una muestra de
10 municipios (dos del Chocó, tres grandes, tres pequeños y los dos sin
edificios de Google que tienen más edificios en Microsoft): cruza los
centroides con los límites con \campo{sjoin(predicate="within")} en el plano y
compara los conteos.

\tabla{validacion_asignacion}

En las 20 comparaciones la diferencia fue cero, sin edificios exclusivos de
ninguno de los dos métodos. La asignación de MongoDB coincide con una
implementación independiente.

\subsection{Agregación en \texorpdfstring{\campo{resultados}}{resultados}}
\label{sec:agregacion}

Cada combinación de fuente y umbral de confianza es una sola pipeline que corre
dentro de MongoDB y escribe directamente en \campo{resultados}:

\begin{lstlisting}[language=Python]
[{"$match": {"cod_mpio": {"$ne": None},
             "confianza": {"$gte": u}}},     # solo Google
 {"$group": {"_id": "$cod_mpio", "n": {"$sum": 1},
             "area": {"$sum": "$area_m2"},
             "media": {"$avg": "$area_m2"}}},
 {"$lookup": {"from": "municipios", "localField": "_id",
              "foreignField": "_id", "as": "m"}},
 {"$unwind": "$m"},
 {"$project": {"_id": {"cod_mpio": "$_id", "fuente_id": f, "umbral": u},
   "n_edificios": "$n", "area_total_m2": "$area", "area_media_m2": "$media",
   "densidad_km2": {"$divide": ["$n", "$m.area_km2"]},
   "pct_techos": {"$multiply": [100, {"$divide":
       ["$area", {"$multiply": ["$m.area_km2", 1e6]}]}]},
   "parametros": {...}, "fecha_calculo": "$$NOW"}},
 {"$merge": {"into": "resultados", "on": "_id",
             "whenMatched": "replace"}}]
\end{lstlisting}

Una segunda pipeline, sobre \campo{municipios}, completa con ceros los
municipios sin edificios (\campo{\$merge} con
\campo{whenMatched: "keepExisting"}), así siempre hay 170 documentos por fuente
y umbral. Antes de escribir se borran los documentos de esa combinación, de
modo que correr el paso dos veces deja exactamente lo mismo. Los umbrales
(0,65, 0,70 y 0,75) se leen de \campo{config.yaml}, y cada documento guarda en
\campo{parametros} el umbral, la regla de asignación, el CRS de las áreas y el
commit del código que lo produjo. Cada pipeline tarda entre 2 y 5\,s.

\subsection{Verificación de los resultados}
\label{sec:verificacion}

Después de escribir, el script comprueba para cada combinación que haya 170
documentos y que la suma de \campo{n\_edificios} sea igual al número de
edificios asignados que pasan el filtro; si algo falla, termina con error. Las
cuatro combinaciones pasaron:

\begin{itemize}
  \item Google 0,65: 2\,694\,785, el total de la colección menos el edificio
        sin municipio.
  \item Google 0,70 y 0,75: 2\,263\,445 y 1\,751\,076, iguales al conteo
        directo con \campo{cod\_mpio} asignado y confianza mayor o igual al
        umbral.
  \item Microsoft: 1\,761\,048, el total de la colección.
\end{itemize}

\subsection{Indicadores derivados para comparar y ordenar}
\label{sec:indicadores}

Para comparar las fuentes y ordenar los municipios se usan tres indicadores,
calculados sobre el escenario base (Google 0,65):

\begin{itemize}
  \item \textbf{Índice de diferencia}
        $D = (A_G - A_M)/(A_G + A_M)$, entre $-1$ (solo Microsoft tiene
        edificios) y $+1$ (solo Google). Se considera que las fuentes
        coinciden si $|D| \leq 0{,}1$ y que están en fuerte desacuerdo si
        $|D| > 0{,}3$.
  \item \textbf{Área conservadora} $\min(A_G, A_M)$: el área de techos que
        las dos fuentes respaldan a la vez. Es el criterio del ranking, para
        que un municipio no suba solo porque una fuente lo sobreestima.
  \item \textbf{Porcentaje de techos}: fracción de la superficie municipal
        cubierta por techos (promedio de las dos fuentes), que indica si el
        área está concentrada en un casco urbano.
\end{itemize}

% =====================================================================
% 8. Resultados
% =====================================================================
\section{Resultados}
\label{sec:resultados}

Salvo que se diga otra cosa, los resultados usan el \textbf{escenario base}:
Google con confianza $\geq 0{,}65$ y Microsoft completo. Las tablas por
municipio están en \ruta{data/processed/resultados_municipio.csv} (680 filas:
170 municipios por cuatro combinaciones de fuente y umbral), y el ranking
completo en \ruta{data/processed/ranking_municipios.csv}.

\subsection{Totales y comparación entre fuentes}

Las dos fuentes casi coinciden en área total (222,17 y
221,13\,km\textsuperscript{2}) aunque Google tiene 1,53 veces más edificios,
porque divide la misma superficie en huellas más pequeñas. La coincidencia
también se da municipio por municipio: la correlación entre el área de las dos
fuentes es de 0,99 (Pearson) y 0,94 (Spearman) sobre los 170 municipios.

\fig[0.85\textwidth]{res_dispersion_municipio.png}{0.4}{Área total por
municipio PDET: Google frente a Microsoft (escala logarítmica). La recta
punteada es 1:1. Los cinco municipios sin edificios de Google no aparecen.}{fig:res-dispersion}

La mayoría de los municipios cae cerca de la recta 1:1. Los que se separan son
pocos, y los más alejados quedan por debajo, es decir, con menos área en
Google: Sipí, López de Micay, Condoto, Carmen del Darién y Ricaurte, todos del
Pacífico.

\subsection{Resultados por subregión}

\tabla{tabla_subregion}

La diferencia entre fuentes no es uniforme. Google tiene más área en el Caribe
(Sierra Nevada - Perijá, Urabá Antioqueño, Montes de María, con $D$ entre
$+0{,}07$ y $+0{,}09$) y Microsoft en el Pacífico y el sur (Chocó, Sur del
Tolima, Pacífico y Frontera Nariñense, con $D$ entre $-0{,}21$ y $-0{,}33$). El
total casi coincide porque los excesos y los faltantes se compensan.

\subsection{Sensibilidad al umbral de confianza}
\label{sec:sensibilidad}

\tabla{sensibilidad}

\fig[0.9\textwidth]{res_sensibilidad_umbral.png}{0.3}{Edificios y área total de
Google según el umbral de confianza, con Microsoft como referencia.}{fig:res-sensibilidad}

Subir el umbral quita muchos más edificios que área: con 0,70 se pierde el
16\,\% de los edificios pero solo el 7\,\% del área, y con 0,75 el 35\,\% y el
19\,\%. Los edificios de baja confianza son, sobre todo, huellas pequeñas. Con
0,75 el área de Google (179,6\,km\textsuperscript{2}) queda 19\,\% por debajo de
Microsoft, así que el umbral sí cambia la conclusión sobre qué fuente tiene más
techo. El orden de los municipios, en cambio, es estable: los 15 con más área
con 0,65 siguen en el top 15 con 0,70, y 14 de 15 con 0,75. Por eso 0,65 es el
escenario base y los otros dos se reportan como sensibilidad.

\subsection{Mapas por municipio}
\label{sec:mapas}

Los mapas (\campo{mapas.py}) usan las mismas clases para las dos fuentes
(quintiles calculados sobre los 340 valores juntos, sin ceros), para que el
mismo color signifique lo mismo en los dos paneles. Las áreas se expresan en
hectáreas (1\,ha $=$ 10\,000\,m\textsuperscript{2}).

\fig{mapa_area_fuentes.png}{0.45}{Área de techos por municipio PDET, Google
(izquierda) y Microsoft (derecha), con las mismas clases. En gris con trama,
los 5 municipios sin edificios en Google.}{fig:mapa-fuentes}

Los dos mapas cuentan la misma historia en casi todo el territorio. El área se
concentra en los municipios con cabeceras urbanas grandes (Santa Marta y
Valledupar superan las 1\,300\,ha con las dos fuentes) y en la Sierra Nevada -
Perijá, Urabá, el norte del Cauca y el piedemonte caqueteño. Los municipios
extensos de la Amazonía y la Orinoquía ocupan buena parte del mapa pero tienen
poca área de techos. La diferencia visible está en el Chocó y la costa
pacífica nariñense, más claros en Google.

\fig[0.9\textwidth]{mapa_diferencia.png}{0.45}{Índice de diferencia $D$ entre
Google y Microsoft por municipio. Naranja: más área en Microsoft; azul: más en
Google; gris: coinciden ($|D| \leq 0{,}1$).}{fig:mapa-diferencia}

En 98 de los 170 municipios las fuentes coinciden dentro de $\pm 10$\,\% y en
otros 49 la diferencia es moderada. Los 23 con desacuerdo fuerte
($|D| > 0{,}3$) no están repartidos al azar: 12 son del Chocó y 6 del Pacífico
y Frontera Nariñense, e incluyen los 5 sin edificios en Google y otros como
Sipí, López de Micay, Ricaurte y Carmen del Darién, donde Google detecta menos
del 10\,\% del área de Microsoft. Solo uno tiene mucha más área en Google:
Istmina ($D = +0{,}60$).

\fig[0.85\textwidth]{mapa_pct_techos.png}{0.42}{Porcentaje de la superficie
municipal cubierta por techos (promedio de las dos fuentes). Clases por
quintiles.}{fig:mapa-pct}

El área total favorece a los municipios grandes; normalizada por la superficie,
la mediana es 0,077\,\%. Piendamó - Tunía (1,33\,\%), Santander de Quilichao
(0,78\,\%), Apartadó (0,65\,\%) y Santa Marta (0,64\,\%) encabezan, y el norte
del Cauca es la zona más densa.

\subsection{Ranking de municipios candidatos}
\label{sec:ranking}

El ranking (\campo{ranking.py}) ordena por área conservadora,
$\min(A_G, A_M)$, y marca los municipios con fuentes en desacuerdo.

\tabla{ranking}

\fig[0.95\textwidth]{ranking_top15.png}{0.4}{Área de techos de los 15
municipios candidatos según cada fuente, ordenados por el área
conservadora.}{fig:ranking}

Los 15 municipios suman 7\,775\,ha de área conservadora, el 39\,\% del total de
los 170 (20\,156\,ha), y ninguno tiene las fuentes en desacuerdo: el criterio
deja fuera, sin necesidad de una regla adicional, a los municipios donde los
datos no son confiables. El ranking es estable: 13 de los 15 coinciden con el
top 15 que saldría de usar solo Google o solo Microsoft, y el umbral de
confianza no cambia cuáles lo encabezan. Están representadas 11 de las 16
subregiones, con Santa Marta, Valledupar y Ciénaga (Sierra Nevada - Perijá) a
la cabeza.

\subsection{Estimación de potencia instalable}
\label{sec:potencial}

El área de techos no es todavía potencial energético. Para dar una idea de la
escala, \campo{potencial.py} convierte el área conservadora de cada municipio
en potencia pico y energía anual con supuestos explícitos, que se leen de
\campo{config.yaml}:

\begin{equation}
  P = A \cdot f_u \cdot \rho, \qquad E = P \cdot H \cdot 365 \cdot PR
  \label{eq:potencial}
\end{equation}

\begin{itemize}
  \item $f_u = 0{,}40$ (rango 0,30 a 0,50): fracción del techo aprovechable
        para módulos, que descuenta sombras, bordes, orientación,
        inclinación y techos que no soportan la carga.
  \item $\rho = 0{,}20$\,kWp/m\textsuperscript{2}: módulos de cerca del
        20\,\% de eficiencia (1\,kW/m\textsuperscript{2} en condiciones
        estándar).
  \item $H = 4{,}5$\,kWh/m\textsuperscript{2}/día: irradiación global
        promedio de Colombia según el Atlas de Radiación Solar de la UPME y
        el IDEAM~\cite{atlas}; se usa un solo valor para todo el país.
  \item $PR = 0{,}80$: rendimiento global del sistema (inversor, temperatura,
        cableado y suciedad).
\end{itemize}

\tabla{potencial}

Con estos supuestos, los 170 municipios PDET tendrían unos 16\,GWp de potencia
instalable en techos (entre 12 y 20\,GWp según el aprovechamiento) y unos
21\,TWh al año, y el top 15 concentra 6,2\,GWp. Es un \textbf{orden de
magnitud}, no un dimensionamiento: es el potencial técnico bruto, que supone
cubrir todos los techos aprovechables, y la irradiación real varía entre
regiones. La sección~\ref{sec:recomendaciones} propone cómo refinarlo.

% =====================================================================
% 9. Discusión y recomendaciones
% =====================================================================
\section{Discusión y recomendaciones}
\label{sec:discusion}

\subsection{Por qué difieren las fuentes}

Los datos muestran tres causas:

\begin{itemize}
  \item \textbf{Tamaño de las huellas.} Google tiene 1,53 veces más edificios
        pero casi la misma área total, porque divide la misma superficie en
        huellas más pequeñas (82\,m\textsuperscript{2} de media frente a
        126\,m\textsuperscript{2}). Cada modelo recorta de forma distinta los
        techos contiguos. Por eso el \emph{conteo} de edificios no sirve para
        comparar fuentes, y el ranking y la estimación de potencial usan
        \emph{área}.
  \item \textbf{Huecos de detección en el Pacífico.} En el Chocó y la costa
        nariñense Google detecta muy poco o nada. Una causa probable es la
        nubosidad casi permanente de la región, que reduce las imágenes
        útiles para el modelo. Ahí se concentran 18 de los 23 municipios con
        desacuerdo fuerte.
  \item \textbf{Fechas e imágenes distintas.} Microsoft usa imágenes de 2014
        a 2021 y Google una inferencia de 2023; una cabecera que creció en
        esos años tiene más área en la fuente más reciente.
\end{itemize}

La consecuencia práctica es que el total regional es robusto (las fuentes
difieren menos del 0,5\,\% y el umbral lo mueve como mucho un 19\,\%), pero la
cifra de un municipio concreto depende de la fuente, sobre todo en el
Pacífico. Por eso se reporta el área conservadora y, donde las fuentes no
coinciden, conviene reportar un rango y no una sola cifra.

\subsection{Recomendaciones para la UPME}
\label{sec:recomendaciones}

\textbf{Selección de pilotos.}
\begin{enumerate}
  \item Partir del top 15 por área conservadora (tabla~\ref{tab:ranking}):
        concentra el 39\,\% del área de techos de los 170 municipios y ninguno
        tiene las fuentes en desacuerdo.
  \item Dentro del top 15, empezar por los municipios con mayor porcentaje de
        techos, donde el área está concentrada en un casco urbano y no
        dispersa en veredas, lo que facilita la conexión y la operación:
        \textbf{Santa Marta}, \textbf{Santander de Quilichao} y
        \textbf{Apartadó} están en el top 15 y entre los cuatro con mayor
        porcentaje de techos.
  \item Buscar diversidad territorial: el top 15 cubre 11 de las 16
        subregiones, así que un conjunto pequeño de pilotos puede
        representar regiones distintas.
\end{enumerate}

\textbf{Municipios del Pacífico.} No descartarlos por su área baja: en los 23
municipios con desacuerdo fuerte el resultado depende más de la fuente que del
municipio. Se recomienda verificarlos antes de decidir, con imágenes
satelitales recientes o de alta resolución, con una tercera fuente (por
ejemplo GlobalBuildingAtlas, que además trae alturas) o en campo, empezando
por los cinco municipios sin edificios en Google.

\textbf{Refinar la estimación antes de dimensionar.}
\begin{enumerate}
  \item Cruzar el área con la irradiación por municipio del Atlas de
        Radiación Solar~\cite{atlas} en lugar del promedio nacional.
  \item Estimar el área útil por techo (sombras, orientación, inclinación,
        material) con imágenes o modelos de elevación en los municipios
        candidatos; el factor de aprovechamiento de 0,40 es un supuesto.
  \item Cruzar con la demanda eléctrica y la red de cada municipio, para
        dimensionar según el consumo y la capacidad de conexión, y no solo
        según el techo disponible.
\end{enumerate}

\textbf{Continuidad de la plataforma de datos.}
\begin{itemize}
  \item Las fuentes se actualizan: el flujo permite recargar una versión
        nueva de cualquiera de ellas con \campo{make edificios} y recalcular
        todo con \campo{make analisis}, conservando la trazabilidad en
        \campo{fuentes} y \campo{bitacora\_cargas}.
  \item Para extender el análisis a todo el país (unos 16,5 millones de
        edificios solo en Google), el diseño no cambia: la asignación ya
        corre municipio por municipio con índice, y MongoDB permite repartir
        las colecciones de edificios con \emph{sharding} por
        \campo{cod\_mpio}.
\end{itemize}

\subsection{Limitaciones}
\label{sec:limitaciones}

\begin{itemize}
  \item El área es la huella en planta: es un límite superior del área útil
        para paneles.
  \item Las huellas vienen de modelos de aprendizaje automático sobre
        imágenes de fechas distintas; no son un catastro y tienen errores de
        omisión y de partición de techos.
  \item La estimación de potencia usa una irradiación nacional y un factor de
        aprovechamiento fijo.
  \item Un centroide justo sobre un límite puede quedar fuera por la
        diferencia entre aristas esféricas y planas (un edificio de 4,5
        millones en este proyecto).
  \item La validación independiente cubre una muestra de 10 municipios, no
        los 170, y el proyecto no tiene pruebas automáticas; la verificación
        se apoya en controles de conteo en cada etapa.
\end{itemize}

% =====================================================================
% 10. Reproducción
% =====================================================================
\section{Reproducción del flujo}
\label{sec:reproduccion}

\subsection{Pasos}

Todo el proyecto se reproduce desde una base vacía con:

\begin{lstlisting}
git clone https://github.com/SamuOsorio/Proyecto-DBA.git && cd Proyecto-DBA
cp .env.example .env
python3 -m venv .venv && source .venv/bin/activate
pip install -r requirements.txt
# copiar el ZIP del MGN 2025 (nivel municipio) en data/raw/mgn_2025/MGN2025.zip
make all        # = make up municipios edificios analisis
\end{lstlisting}

\campo{make analisis} corre \campo{src/analisis/run.py}, que ejecuta los pasos
en orden como \campo{python -m src.analisis.<paso>}; si uno falla, se detiene
e indica cómo retomar con \campo{--desde <paso>}.
\campo{--solo-figuras} rehace gráficas, mapas, ranking y potencial a partir de
los CSV y el GeoJSON versionados, sin necesidad de MongoDB.

\begin{table}[H]
  \centering\small
  \caption{Pasos de \campo{make analisis}, en orden.}
  \label{tab:pasos}
  \begin{tabularx}{\textwidth}{@{}l L@{}}
    \toprule
    \textbf{Paso} & \textbf{Qué produce} \\
    \midrule
    \campo{asignacion} & \campo{cod\_mpio} en los edificios; registro en \campo{bitacora\_cargas} \\
    \campo{evidencia} & \campo{explain()} con IXSCAN y con COLLSCAN \\
    \campo{validacion} & Comparación con GeoPandas en 10 municipios \\
    \campo{agregacion} & \campo{resultados}: 170 documentos por fuente y umbral, verificados \\
    \campo{sensibilidad} & Tabla de sensibilidad al umbral de Google \\
    \campo{exportar} & CSV por municipio y subregión; tabla por subregión \\
    \campo{figuras\_eda} & Las figuras del EDA \\
    \campo{graficas} & Dispersión por municipio y gráfica de sensibilidad \\
    \campo{mapas} & Mapas por fuente, de diferencia y de \% de techos \\
    \campo{ranking} & Top 15, tabla y gráfica; ranking completo en CSV \\
    \campo{potencial} & Potencia instalable y energía por municipio \\
    \bottomrule
  \end{tabularx}
\end{table}

\subsection{Corrida desde base vacía}

Para este informe se borró la base \campo{upme\_solar} y se corrió
\campo{make all}. Terminó sin errores en unos 14 minutos, incluyendo el
prefiltro y la limpieza de Google desde los archivos crudos, y reprodujo
exactamente las mismas cifras: 170 municipios, 2\,694\,786 y 1\,761\,048
edificios, el mismo edificio sin municipio y 680 documentos en
\campo{resultados}. Los CSV regenerados solo difieren en la última cifra
decimal, por redondeo de punto flotante.

\fig[0.75\textwidth]{cap_make_all.png}{0.32}{Corrida de \campo{make all}
desde una base vacía: resumen de tiempos de la etapa de análisis
(\campo{run.py}).}{fig:make-all}

\subsection{Parámetros y entorno}

Los parámetros no están escritos en el código: los umbrales de confianza, la
regla de asignación, los CRS, el tamaño de los lotes, los supuestos de la
estimación de potencia y las rutas de salida se leen de \campo{config.yaml}.
La URI de la base vive en \campo{.env}, fuera del repositorio. Cada documento
de \campo{resultados} guarda el commit del código que lo produjo.

Entorno: MongoDB 8 en Docker Compose y Python 3.12 o superior con las
versiones de \campo{requirements.txt}. El flujo se corrió en macOS y en Linux,
con las mismas cifras.

\subsection{Lecciones del proceso}

\begin{itemize}
  \item Un validador \campo{\$jsonSchema} copiado tal cual del diccionario
        bloquea el flujo si exige campos que se calculan en una etapa
        posterior; hubo que declararlos opcionales o nulos.
  \item Que una geometría sea válida en el plano (GEOS) no garantiza que lo
        sea en la esfera (S2): crear el índice \campo{2dsphere} antes de
        cargar convirtió un fallo de todo el índice en el rechazo de unos
        pocos documentos reparables.
  \item Separar cada tarea en su propio archivo, comunicadas por archivos
        intermedios y colecciones, permitió que cuatro personas trabajaran en
        paralelo, a cambio de coordinar el orden de los pasos en
        \campo{run.py}.
  \item Guardar los CSV de resultados en el repositorio permitió rehacer
        mapas y figuras sin tener la base cargada, lo que aceleró la revisión
        de las gráficas.
\end{itemize}

% =====================================================================
% 11. Conclusiones
% =====================================================================
\section{Conclusiones}
\label{sec:conclusiones}

El proyecto entrega a la UPME una metodología reproducible, construida sobre
MongoDB, para contar las edificaciones y estimar el área de techos de los 170
municipios PDET con dos fuentes abiertas. El flujo va de los límites oficiales
del DANE y las huellas de Google y Microsoft hasta una colección
\campo{resultados} con 680 documentos verificados contra los conteos de las
colecciones, y de ahí a tablas, mapas, un ranking y una estimación de potencia,
todo con un solo comando y desde una base vacía.

Las decisiones de base de datos sostienen el resultado: los índices
\campo{2dsphere} hacen que la asignación de 4,5 millones de edificios use
\campo{IXSCAN} y no recorra las colecciones; la regla del centroide evita
contar dos veces un techo en el límite; la agregación corre dentro de la base;
los validadores \campo{\$jsonSchema} y el índice único sobre \campo{geom\_hash}
garantizan la integridad; y \campo{fuentes} y \campo{bitacora\_cargas} dejan
trazado de dónde viene cada dato y con qué parámetros se calculó. La asignación
coincidió al 100\,\% con una implementación independiente en GeoPandas.

En cuanto a los resultados, las dos fuentes coinciden en el total (222 y
221\,km\textsuperscript{2}) y en el orden de los municipios, pero no en el
Pacífico, donde Google tiene huecos de detección. El área se concentra: 15
municipios, encabezados por Santa Marta, Valledupar y Buenaventura, suman el
39\,\% del área conservadora, con una potencia instalable del orden de
6\,GWp. Esos municipios son los candidatos naturales para los pilotos; los del
Pacífico requieren verificación antes de decidir, y cualquier dimensionamiento
debe refinar el área útil y cruzarla con la irradiación y la demanda local.


% ------------------------------- Referencias -------------------------
\begin{thebibliography}{9}

\bibitem{enunciado}
A.~O. Calderón Romero, \emph{Database Administration -- Final Project},
Pontificia Universidad Javeriana, 2026.

\bibitem{pdet}
Agencia de Renovación del Territorio, \emph{¿Qué son los PDET?} y listado
\texttt{MunicipiosPDET.xlsx}.
\url{https://www.renovacionterritorio.gov.co/central-pdet/que-son-los-pdet}

\bibitem{mgn}
DANE, \emph{Marco Geoestadístico Nacional -- Descarga de datos
geoestadísticos}, versión MGN2025.
\url{https://geoportal.dane.gov.co/servicios/descarga-y-metadatos/datos-geoestadisticos/?cod=111}

\bibitem{google}
W.~Sirko et al., \emph{Open Buildings V3}, Google Research.
\url{https://sites.research.google/gr/open-buildings/}

\bibitem{microsoft}
Microsoft, \emph{Global ML Building Footprints}.
\url{https://github.com/microsoft/GlobalMLBuildingFootprints};
\url{https://planetarycomputer.microsoft.com/dataset/ms-buildings}

\bibitem{mongodb}
MongoDB, Inc., \emph{Geospatial Queries} y \emph{2dsphere Indexes},
documentación de MongoDB 8.
\url{https://www.mongodb.com/docs/manual/geospatial-queries/}

\bibitem{atlas}
UPME e IDEAM, \emph{Atlas de Radiación Solar, Ultravioleta y Ozono de
Colombia}.
\url{http://atlas.ideam.gov.co/visorAtlasRadiacion.html}

\end{thebibliography}

\appendix
% =====================================================================
% Anexo: evidencias de ejecución
% =====================================================================
\section{Evidencias de ejecución}
\label{sec:anexo}

Capturas de las ejecuciones reales contra MongoDB, tomadas durante el
proyecto, en el orden del flujo. Los informes de cada entrega
(\ruta{Entrega 1/} a \ruta{Entrega 4/} en el repositorio) tienen el detalle de
cada paso.

\subsection{Límites PDET}

\fig{cap_cruce.png}{0.25}{\campo{cruce.py}: localiza la capa municipal del MGN
y cruzan los 170 municipios PDET.}{fig:a-cruce}
\fig{cap_geometria.png}{0.3}{\campo{geometria.py} y \campo{exportar.py}: 170
municipios revisados, ninguno inválido, y el GeoJSON exportado.}{fig:a-geometria}
\fig{cap_schema.png}{0.2}{\campo{schema.py}: crea \campo{municipios} y
\campo{fuentes}; al volver a correrlo no modifica las existentes.}{fig:a-schema}
\fig{cap_validacion.png}{0.25}{Un documento sin los campos obligatorios es
rechazado por el validador \campo{\$jsonSchema} (error 121).}{fig:a-validacion}
\fig{cap_carga.png}{0.2}{\campo{carga.py}: los 170 municipios en la colección.}{fig:a-carga}
\fig{cap_indices.png}{0.2}{\campo{indices.py}: \campo{geometry\_2dsphere} y
\campo{subregion\_pdet.id\_1}.}{fig:a-indices}
\fig{cap_evidencia.png}{0.25}{\campo{evidencia.py}: el plan usa \campo{IXSCAN} y
\campo{\$geoIntersects} devuelve el municipio correcto.}{fig:a-evidencia}

\subsection{Google Open Buildings}

\fig{cap_google_listar.png}{0.25}{Celdas S2 que tocan Colombia y su tamaño
(\texttt{--listar}).}{fig:a-g-listar}
\fig{cap_google_descarga.png}{0.25}{Descarga o reutilización de las celdas y
SHA-256 de la fuente.}{fig:a-g-descarga}
\fig{cap_google_fuentes.png}{0.3}{Documento de Google en \campo{fuentes}.}{fig:a-g-fuentes}
\fig{cap_run_google.png}{0.25}{Tabla de control de Google en \campo{make
edificios}.}{fig:a-g-run}

\subsection{Microsoft Building Footprints}

\fig{cap_ms_indice.png}{0.2}{Índice de Microsoft: 232 particiones de
Colombia.}{fig:a-ms-indice}
\fig{cap_ms_descarga.png}{0.2}{Descarga: 125 de 232 particiones tocan PDET y
SHA-256 global.}{fig:a-ms-descarga}
\fig{cap_ms_manifest.png}{0.22}{\campo{manifest.csv}: tamaño y SHA-256 de cada
archivo.}{fig:a-ms-manifest}
\fig{cap_ms_parquet.png}{0.25}{Verificación del Parquet de Microsoft con
DuckDB.}{fig:a-ms-parquet}
\fig{cap_ms_geoparquet.png}{0.25}{Lectura de los Parquet de Microsoft y de
Google con GeoPandas.}{fig:a-ms-geoparquet}
\fig{cap_ms_fuentes.png}{0.2}{Documento de Microsoft en \campo{fuentes}.}{fig:a-ms-fuentes}
\fig{cap_ms_bitacora.png}{0.35}{Entradas de \campo{bitacora\_cargas} de
Microsoft.}{fig:a-ms-bitacora}
\fig{cap_run_microsoft.png}{0.18}{Tabla de control de Microsoft en \campo{make
edificios}.}{fig:a-ms-run}

\subsection{Carga e índices de edificios}

\fig{cap_schema_edificios.png}{0.3}{\campo{schema.py}: colecciones de edificios
con validación \campo{\$jsonSchema}.}{fig:a-schema-edif}
\fig{cap_validacion_edificios.png}{0.3}{Un edificio sin campos obligatorios es
rechazado (error 121).}{fig:a-validacion-edif}
\fig{cap_carga_edificios.png}{0.35}{\campo{carga.py} para las dos fuentes.}{fig:a-carga-edif}
\fig{cap_indices_edificios.png}{0.3}{\campo{indices.py} con el tiempo de cada
índice.}{fig:a-indices-edif}
\fig{cap_evidencia_edificios.png}{0.35}{\campo{evidencia.py}: planes con
\campo{IXSCAN} y tamaños de las colecciones.}{fig:a-evidencia-edif}
\fig{cap_run_edificios.png}{0.3}{Final de \campo{make edificios}: carga, índices y
evidencia.}{fig:a-run-edif}

\subsection{Análisis}

\fig{cap_asignacion.png}{0.35}{Salida de la asignación por centroide.}{fig:a-asignacion}


\end{document}
